\section{Method}
\label{sec:method}

We first describe the agent and its editable harness (\S\ref{sec:agent}), then our port of the
Continual Harness refiner to local models (\S\ref{sec:refiner}), the five techniques we add on top
of it (\S\ref{sec:techniques}), and the baselines (\S\ref{sec:baselines}).

\subsection{Agent and Harness}
\label{sec:agent}

\paragraph{Agent loop.}
The agent is a standard tool-calling loop in the style of ReAct~\cite{yao2023react}. At each step
the model receives the conversation so far and either calls one or more tools or replies without a
tool call, which ends the task. A task also ends after a maximum of 15 steps. The model has three
base tools: \texttt{run\_shell} (run a shell command in the task folder), \texttt{read\_file} and
\texttt{write\_file}. Every task runs in its own fresh temporary folder. Shell commands and skill code
run inside an operating-system sandbox that allows writing only inside the task folder, blocks
reading the user's files, has no network access, and stops any command after 30 seconds.

\paragraph{Harness.}
Following Continual Harness~\cite{karten2026continual}, the \emph{harness} is everything around the
model weights that shapes its behaviour, and it has four editable parts:
\begin{itemize}
  \item \textbf{Strategy}: a free-text ``base prompt'' added to a fixed system prompt. The fixed
        part (tools and hard rules) is never edited. The initial strategy is one sentence:
        \emph{``Do the task step by step. Check your work (for example run the tests) before you finish.''}
  \item \textbf{Memory}: a list of entries (title, content, importance from 1 to 5) shown in the
        system prompt.
  \item \textbf{Skills}: named pieces of Python code that the agent can run with a
        \texttt{run\_skill} tool. The code runs in the task folder, inside the sandbox.
  \item \textbf{Sub-agents}: named helper agents, each with its own instructions, an allowed list
        of tools and a turn limit, that the agent can call with a \texttt{call\_subagent} tool.
\end{itemize}
The \texttt{run\_skill} and \texttt{call\_subagent} tools are only offered once the harness contains
at least one skill or sub-agent. As a result, before the first edit, the self-improving agent is
exactly the static baseline agent. To keep prompt size, memory use and run time bounded on small
hardware, every part has a size cap: strategy 4{,}000 characters; 20 memory entries of up to 500
characters; 10 skills of up to 3{,}000 characters of code; 5 sub-agents with at most 8 turns each.

\subsection{Continual Harness Refiner (C1)}
\label{sec:refiner}

The original Continual Harness code supports only cloud model APIs and is tied to game emulators,
so we re-implemented its refiner (\texttt{HarnessEvolver}) for local models served by
Ollama~\cite{ollama} and for our tasks, keeping its design as close as possible.

\paragraph{What we keep.}
\begin{itemize}
  \item \textbf{Schedule.} No edits before step 25 (warm-up); then an edit every 25 steps up to
        step 200, and every 100 steps after that. Steps are counted across tasks.
  \item \textbf{Window.} The refiner reads the last 25 steps of the agent's trajectory (the last 20
        for the memory pass).
  \item \textbf{Four passes,} run in order, each in its own error handler so that one failing pass
        does not stop the others: (1) rewrite the strategy; (2) create, update or retire sub-agents;
        (3) add or update skills; (4) add or update memory entries. The strategy pass returns
        markdown; the other passes return JSON, parsed the same way as in the original.
  \item \textbf{Same model.} The model that does the tasks also writes the edits.
  \item \textbf{No acceptance test.} Every edit is kept.
\end{itemize}

\paragraph{What we change.}
\begin{itemize}
  \item Edits happen \emph{between} tasks: if a task crosses one of the scheduled step counts, the
        refiner runs after that task ends.
  \item The trajectory shown to the refiner marks the end of each task with its result (PASS or
        FAIL). This replaces game progress, which the original refiner sees.
  \item Prompts use task words instead of game words, skill code runs as a Python file in the
        sandbox instead of calling a game API, and sub-agents have fewer fields.
  \item A skill's code is saved when the skill is added (the original saves code only on update).
  \item The size caps above, and full logging of every edit (before and after).
\end{itemize}
We call this port, with nothing added, condition \textbf{C1}.

\subsection{Added Techniques (C2--C6)}
\label{sec:techniques}

Each technique targets a weakness of C1 that we expect to matter most for small models. We test
each one alone on top of C1, and all of them together (C8).

\paragraph{Separate refiner.}
A larger local model writes the edits while the small model still does the tasks. This tests
whether the bottleneck is the quality of the edits. For both workers we use
\texttt{gemma4:e4b} as the refiner.

\paragraph{Regression gate.}
After the refiner edits the harness, the gate reruns the most recent train tasks that ran with the
old harness since the last edit (at most 3), with the same seeds, using the new harness. The edits
are kept only if the new harness passes at least as many of these tasks as the old one did;
otherwise the harness is reset to its state before the edit. This is the gated-edit idea of
Self-Harness~\cite{zhang2026selfharness}, scaled down to a small task budget.

\paragraph{Memory database.}
In C1, every memory entry goes into every prompt, and memory is capped at 20 entries. With the
memory database, the refiner can store up to 200 entries, and at the start of each task only the 5
entries most similar to the task instruction are put into the prompt. Similarity is the cosine
similarity of \texttt{nomic-embed-text} embeddings~\cite{nussbaum2024nomic}.

\paragraph{Verified skills.}
In C1, any skill the refiner writes is kept. With verified skills, every new or rewritten skill is
first run once on a fresh copy of the folder of the most recent train task, and it is kept only if
it runs without an error and returns a result. This is a weaker form of the code verification used
by skill-library agents such as Voyager~\cite{wang2023voyager}.

\paragraph{Structured errors.}
In C1, a failed tool call returns the raw output (exit code, standard output and standard error).
With structured errors, a failed call returns a short header before the raw output: one of 14 fixed
error types (for example \texttt{file\_not\_found}, \texttt{syntax\_error},
\texttt{test\_failed}), the key line of the message (for example the last line of a Python
traceback), and a fixed, general hint for that type. Hints never mention a specific task.
Successful calls are unchanged.

\paragraph{Full system (C8).}
C1 with all five techniques together.

\subsection{Baselines}
\label{sec:baselines}

\paragraph{Static harness (C0).}
The same agent, tools and sandbox with a fixed system prompt and no learning. This is the main
reference point: a self-improving harness is only useful if it beats C0.

\paragraph{Reflexion-style memory.}
Following Reflexion~\cite{shinn2023reflexion}, adapted to a sequence of different tasks: after each
\emph{failed} train task, the same model reads the task and its own steps and writes a short
self-reflection (up to 500 characters) that diagnoses the failure and gives a plan for similar
tasks. The last 3 reflections are shown as memory in the system prompt. The strategy, skills and
sub-agents never change.

\paragraph{Gated edits.}
The regression gate condition above serves as our Self-Harness-style baseline.
